% Generated draft: requires mathematical and source review.
\documentclass[11pt]{article}
\usepackage[margin=25mm]{geometry}
\usepackage{fontspec}
\setmainfont{Latin Modern Roman}
\usepackage{amsmath,amssymb,mathtools,unicode-math}
\setmathfont{Latin Modern Math}
\usepackage{enumitem,needspace,xurl,hyperref}
\setlength{\parindent}{0pt}
\setlength{\parskip}{5pt}
\setlength{\emergencystretch}{3em}
\newcommand{\N}{\mathbb{N}}
\newcommand{\Z}{\mathbb{Z}}
\newcommand{\Q}{\mathbb{Q}}
\newcommand{\R}{\mathbb{R}}
\newcommand{\C}{\mathbb{C}}
\newcommand{\F}{\mathbb{F}}
\newcommand{\A}{\mathbb{A}}
\newcommand{\PP}{\mathbb{P}}
\newcommand{\E}{\mathbb{E}}
\newcommand{\eps}{\varepsilon}
\newcommand{\dd}{\,\mathrm{d}}
\begin{document}
\section*{UnsolvedMath ID: 11000198}
\textbf{Source problem code:} catalogue entry
\subsection{Definitions and mathematical statement}
Retain the source's surface \(\Sigma\), group \(\pi\), higher-rank representation space, and energy function \(E_\rho\). In the conventional setting, \(\pi=\pi_1(\Sigma)\), \(\rho:\pi\to G\) is a representation into a Lie group with a symmetric-space target, and \(E_\rho\) is an energy function on Teichm\"uller space, defined through equivariant harmonic maps under appropriate existence hypotheses. Those data and hypotheses are not supplied in the excerpt.

\emph{Problem 3.6.} Find a replacement for convex cocompactness in higher rank that includes the preceding examples of proper \(\operatorname{Mod}_\Sigma\)-actions and for which \(E_\rho\) is proper. A candidate must specify the higher-rank group, the representation condition, the mapping-class-invariant domain and action under consideration, and the definition and domain of \(E_\rho\). For an energy \(E_\rho:\operatorname{Teich}(\Sigma)\to[0,\infty)\), properness means compact preimages of compact sets; in the usual continuous setting it requires compact bounded-energy sublevel sets. The source's preceding examples, which the condition must include, are absent and cannot be replaced by an arbitrarily chosen family.

The source motivates the requirement by reporting that, for a discrete embedding \(\pi\hookrightarrow\operatorname{PSL}_2(\C)\), the energy function is proper if and only if the representation is quasi-Fuchsian. It attributes this to Bonahon--Thurston geometric tameness and then-recent extensions, citing [53]. The precise surface, parabolic, and energy hypotheses of that equivalence are omitted and must be checked. This rank-one statement is contextual evidence, not a supplied higher-rank criterion, and ``recent'' is relative to the source.
\subsection{Short English statement}
Develop a higher-rank analogue of convex cocompactness that covers the source's proper mapping-class actions and guarantees proper energy functions. The examples and energy setup are omitted; the reported quasi-Fuchsian equivalence supplies motivation rather than a higher-rank answer.
\subsection{Sources}
\begin{itemize}
\item Local UnsolvedMath registry, permanent ID 11000198, source code catalogue entry, Problem 3.6, PDF page 225. No individual proposer is supplied; page distinguishes repeated labels.
\item Farb, \emph{Problems on Mapping Class Groups and Related Topics}. Supplied URL: \url{https://www.math.uchicago.edu/~farb/papers/mcgbook.pdf}; recorded access 2026-07-29; extraction pdf-font-structure. The edited volume contains 329 explicitly typeset headings across chapters.
\item Bonahon--Thurston and their geometric-tameness extensions are credited in the statement; reference [53] is supplied for the energy/quasi-Fuchsian equivalence but has no full bibliographic entry.
\item Supplied status open and OPEN-TRIAGE. The summary describes a broad higher-rank Anosov/convex-cocompact analogue program and says it is active and unsettled, without a precise result answering the stated energy requirement. Current primary-literature status and redistribution terms remain NEEDS\_REVIEW.
\end{itemize}
\end{document}
